% Pillar-5 (MIN-REPORT-RL) central standard: the 7-item minimum reportable stack
\section{The Seven-Item Minimum Reportable Stack}
\label{sec:p5-stack}
A reporting standard survives only if it is short, checkable, and every item is
load-bearing. We therefore admit an item to the minimum reportable stack only
if we can point to a \emph{documented lever}: evidence that varying the item,
with the rest held fixed, moved or flipped a reported outcome or its
telemetry. \tableref{tab:p5-minreport} is the standard; the subsections justify
each row. Items 1--3 specify what objective was optimized and on what
substrate; items 4--5 specify the signal geometry of group-relative training;
items 6--7 specify what the headline number can legitimately claim.

\begin{table}[t]
\centering
\caption{The seven-item minimum reportable stack. Each item earns its place
via a documented lever---evidence that the item alone can move or flip a
reported result. Provenance strata as defined in \secref{sec:p5-results}.}
\label{tab:p5-minreport}
\footnotesize
\resizebox{\textwidth}{!}{%
\begin{tabular}{@{}cp{2.4cm}p{4.4cm}p{5.6cm}@{}}
\toprule
\# & Item & What must be reported & Documented lever \\
\midrule
1 & Loss form & Ratio level (token/sequence), clip values and asymmetry,
    token mask, advantage normalization, dynamic-sampling status &
    Same label, opposite ZVF telemetry across loss forms
    (0.00 vs.\ 0.58, \secref{sec:p5-exhibit-label}); normalization biases
    \citep{tong2025drgrpo} \\
2 & Reference policy \& KL & Reference snapshot, KL coefficient and estimator,
    or explicit ``no KL'' &
    Distinguishes GRPO-family members that are otherwise conflated
    \citep{shao2024deepseekmath, yu2025dapo}; silently alters the objective \\
3 & Sampler / backend / precision & Training backend, sampling engine,
    precision, decoding parameters &
    $17\times$ outcome flip from a backend swap alone
    (\secref{sec:p5-exhibit-backend}) \\
4 & Per-step ZVF/GU trajectory & Per-step zero-variance fraction and gradient
    utilization $\mathrm{GU} = 1-\mathrm{ZVF}$ &
    ZVF is U-shaped in accuracy (\tableref{tab:p5-ushape}): a mean hides
    whether signal starved from mastery or incapacity \\
5 & Group-size schedule & $G$ per step, including adaptive escalation rules &
    ZVF falls systematically with $G$ (\tableref{tab:p5-gsize}); theory:
    $\E[\mathrm{ZVF}] = p^G + (1-p)^G$ at per-prompt accuracy $p$ \\
6 & Held-out split & Split identity, size, and disjointness from the reward
    environment &
    Train reward and held-out accuracy dissociate; reward-side
    overoptimization is well documented \citep{gao2023reward} \\
7 & Decontamination \& parser probe & Contamination check performed;
    sensitivity of reward/telemetry to parser perturbation &
    Benchmark contamination shifts headline scores \citep{zhang2024gsm1k,
    riddell2024contamination}; micro-jitter $\epsilon \sim U(0,10^{-4})$
    collapses batch ZVF $0.158 \to 0.000$ \\
\bottomrule
\end{tabular}}
\end{table}

\subsection{Item 1: Loss Form}
\label{sec:p5-item1}
The label is not the loss. ``DAPO'' with and without dynamic sampling are
different optimizers with opposite ZVF signatures
(\secref{sec:p5-exhibit-label}); Dr.GRPO exists because GRPO's
std-normalization and length terms bias the update \citep{tong2025drgrpo};
GSPO moves the importance ratio from token to sequence level
\citep{qwen2025gspo}; and the group-relative loss family is close enough to
preference optimization that GRPO admits a DPO reading
\citep{rafailov2024direct, wu2025grpo_dpo}, making the exact contrast structure
part of the objective's identity. The reportable unit is therefore the tuple
(ratio level, clip low/high, token mask, advantage normalization,
dynamic-sampling status), not the name.

\subsection{Item 2: Reference Policy and KL Handling}
\label{sec:p5-item2}
Whether a run regularizes toward a frozen reference, a moving reference, or
nothing at all changes what is being optimized; GRPO as introduced carries a
KL term \citep{shao2024deepseekmath} that many descendants drop or re-estimate
differently. Because the KL term interacts with the clip range in determining
how far the policy can move per step, two runs labeled identically but
differing in KL handling can occupy different stability regimes. This item is
cheap to report (three fields) and its omission is unrecoverable post hoc.

\subsection{Item 3: Sampler, Backend, and Precision}
\label{sec:p5-item3}
Exhibit 1 is the lever: a backend swap alone produced the largest effect in
our entire corpus ($5.0\% \to 84.4\%$). Serving engines differ in kernel
selection, batching, and numerics \citep{kwon2023vllm, zheng2024sglang,
dao2022flashattention}, and in on-policy RL these differences compound: the
sampler's distribution \emph{is} the data distribution. A closed backend is
not disqualifying---but it must be named, versioned where possible, and its
precision and decoding parameters stated, so that a reader knows the result is
conditioned on it.

\subsection{Item 4: The Per-Step ZVF/GU Trajectory}
\label{sec:p5-item4}
A single summary number cannot substitute for the trajectory, because ZVF
aliases opposite regimes. Across a 368-run audit (Stratum B; drawn from a
790-run corpus), mean ZVF is U-shaped in accuracy (\tableref{tab:p5-ushape}):
near-zero-reward runs and near-saturated runs both show ZVF near $0.9$--$1.0$,
while mid-accuracy runs sit near $0.25$. Theory makes the aliasing exact: with
per-prompt accuracy $p$ and binary rewards, the expected zero-variance
indicator is $p^G + (1-p)^G$, symmetric in $p \leftrightarrow 1-p$, and the
usable signal scales as $p(1-p)$ (validated directionally at toy scale,
Stratum C: corr(grad\_norm, $p(1-p)$) $= +0.71$ on a 0.5B model with an open
backward pass---an effect direction, not a publishable effect size). Reporting
the per-step trajectory (and $\mathrm{GU} = 1 - \mathrm{ZVF}$) lets readers
see \emph{which} arm of the U a run occupied; the companion Pillar-2 paper
develops the diagnostic and its failure modes in full.

\begin{table}[t]
\centering
\caption{ZVF is U-shaped in accuracy (368-run audit; internal program
records, June 2026). Mean reward and mean ZVF per model; ZVF is high at both
extremes and lowest mid-scale, so a run's mean ZVF is uninterpretable without
its accuracy regime.}
\label{tab:p5-ushape}
\small
\begin{tabular}{@{}lcc@{}}
\toprule
Model & Mean reward & Mean ZVF \\
\midrule
Llama-3.2-1B        & 0.01 & 0.97 \\
Llama-3.2-3B        & 0.01 & 0.95 \\
Qwen3-32B           & 0.35 & 0.25 \\
Qwen3-8B            & 0.50 & 0.29 \\
Nemotron-30B        & 0.77 & 0.50 \\
Qwen3-4B-Instruct   & 0.86 & 0.87 \\
Qwen3-30B-Instruct  & 0.95 & 0.96 \\
\bottomrule
\end{tabular}
\end{table}

\subsection{Item 5: Group-Size Schedule}
\label{sec:p5-item5}
Group size is a signal-geometry knob, not a throughput knob: larger $G$ lowers
the probability that a group is degenerate (from the same expression
$p^G + (1-p)^G$), and the effect is large in practice
(\tableref{tab:p5-gsize}). Because adaptive schedules now exist---our
open-trainer audit escalated $G$ from 4 to 6 to 8 live on ZVF spikes
(\secref{sec:p5-toolchain})---reporting a single scalar $G$ is no longer
sufficient; the schedule is part of the optimizer. The companion Pillar-3
paper of \ours{} treats group size in depth, including the rollout-budget
trade-off highlighted by 2-GRPO-style analyses \citep{wu2025grpo_dpo,
tan2025scalingrl}.

\begin{table}[t]
\centering
\caption{Mean ZVF by model and group size $G$ (internal program records,
June 2026; ``--'' = not run). Larger $G$ lowers ZVF where accuracy is
intermediate, as theory predicts; near the extremes of the U-shape the effect
compresses.}
\label{tab:p5-gsize}
\small
\begin{tabular}{@{}lccc@{}}
\toprule
Model & $G=4$ & $G=8$ & $G=16$ \\
\midrule
Llama-3.2-1B      & 0.98 & 0.97 & --   \\
Qwen3-4B-Instruct & 0.89 & 0.84 & 0.91 \\
Llama-3.1-8B      & 0.76 & 0.65 & 0.47 \\
Qwen3-8B          & 0.36 & 0.27 & 0.21 \\
Qwen3-32B         & 0.33 & 0.18 & 0.08 \\
\bottomrule
\end{tabular}
\end{table}

\subsection{Item 6: A Held-Out Split Disjoint from the Reward Environment}
\label{sec:p5-item6}
Training reward is measured inside the same apparatus that generates the
learning signal, so it inherits every quirk of that apparatus---parser
leniency, template artifacts, reward misspecification
\citep{gao2023reward, krakovna2020specification}. In the \ours{} group-size
sweep, train-side telemetry moved substantially while held-out accuracy was
flat, a dissociation that would be invisible without a disjoint split. The
requirement is deliberately minimal: name the split, state its size, and
certify it is disjoint from the prompts and the verifier loop used for
training.

\subsection{Item 7: Decontamination and a Parser-Robustness Probe}
\label{sec:p5-item7}
Contamination inflates headline numbers \citep{zhang2024gsm1k,
riddell2024contamination}; that much is standard. The less familiar half of
this item is the parser probe: because rewards are computed by string-level
verifiers, sub-reward perturbations can flip telemetry wholesale. In our
convergent repository experiment, adding micro-jitter
$\epsilon \sim U(0, 10^{-4})$ to rewards collapsed batch ZVF from $0.158$ to
$0.000$ while a magnitude-aware alternative (pairwise contrast density,
$\mathrm{PCD} = \frac{G-1}{G}\,\E[p(1-p)]$) moved by $3\times 10^{-7}$. A
diagnostic that a parser's rounding behavior can zero out is a diagnostic
whose reported value is a stack property. The probe we ask for is one line:
perturb rewards below the verifier's resolution and report which telemetry
survives.

\subsection{Live-Corpus Coupling: How the 7-Item MIN-REPORT Fingerprint Relates to Measured Telemetry}
\label{sec:p5-live-coupling}
The seven items above are the proposed MIN-REPORT standard; the question a
reviewer will ask is whether the resulting manifest \emph{actually predicts}
the measured telemetry on a live corpus. We test this on the
98-cell mega-campaign manifest corpus
(\texttt{experiments/results/mega\_20260704/manifests/}, paired with
\texttt{cells.tsv}; analysed by
\texttt{scripts/p5p8/p5\_manifest\_outcome\_coupling.py}, see
\texttt{experiments/results/p5p8/p5\_manifest\_outcome\_coupling\.*}). On
the live corpus, the 7-item manifest fingerprint carries a total of
$11.4$ bits of cross-cell information; four of the seven items
(\textit{loss\_form}, \textit{ref\_policy\_kl},
\textit{sampler\_backend\_precision}, and the per-cell
\textit{per\_step\_zvf\_path}, which is unique-per-cell and thus
behaves as a cell identifier rather than a stack descriptor) contribute
exactly $0$ bits, while the remaining three stack-descriptor items
(\textit{group\_size\_schedule}, \textit{heldout\_split},
\textit{decontamination\_notes}) carry all $11.4$ bits. The badge is
therefore constant across cells and cannot discriminate seed-pair
reward dispersion within a cell-group (Spearman $\rho = +0.000$, CI
$[+0.000, +0.000]$, $n=48$ pairs); but the manifest \emph{fingerprint as a
whole} is a strong predictor of joint measured telemetry. Across
$2000$ sampled cell-pairs, the Hamming distance on the 7-item manifest
predicts the absolute difference on ZVF (Spearman $\rho = +0.529$,
close-vs-far permutation $p \le 10^{-3}$), on PCD
($\rho = +0.541$, $p \le 10^{-3}$), and on mean-reward
($\rho = +0.251$, $p = 0.003$). The pattern is concentrated entirely in
the three items that carry cross-cell information: per-item permutation
max-$|\Delta\text{mean}|$ tests return $p \le 10^{-3}$ for
\textit{heldout\_split} and \textit{decontamination\_notes} on all three
channels and for \textit{group\_size\_schedule} on ZVF/PCD, while the
four placebo items return NaN (no test possible) or $p = 1.000$. The
operational reading: a manifest-emit pipeline that varies only on
operational metadata makes MIN-REPORT \emph{auditable} but not
\emph{discriminative} on the algorithm axis; the standard's predictive
power is bounded by what the emitters actually vary on. The sharp
operational recommendation that falls out of this measurement is
additive: tag each emit with the set of stack axes it varies on, so
downstream consumers can automatically report the axis-discrimination
ceiling alongside the badge.
